\documentclass[a4paper,12pt]{article}
\usepackage[bahasa]{babel}
\usepackage{graphicx}
\usepackage{multirow}
\usepackage{enumitem}
\usepackage{listings}
\usepackage{adjustbox}
\graphicspath{ {./img/} }
\begin{document}
\title{Laporan Praktikum Statistika Pertemuan 13}
\author{Aldzikri Dwijayanto Prathama 
	\\195410189}
\makeatletter
\begin{titlepage}
	\begin{center}
		{\huge \bfseries \@title }\\[14ex]
		\includegraphics[scale=.8]{logo}\\[4ex]
		{\large \@author}\\[20ex]
		{\large \bfseries {SEKOLAH TINGGI MANAJEMEN INFORMATIKA DAN KOMPUTER
				AKAKOM YOGYAKARTA}}
	\end{center}


\end{titlepage}
\makeatother
\newpage
\tableofcontents
\newpage
\section{Tujuan}
Praktikan mampu melakukan analisis data menggunakan uji rata-rata satu populasi normal dan variansi tidak diketahui

\section{Dasar Teori}
\subsection{Uji hipotesis Mean Populasi Normal}
Ingin diketahui apakah mean ($\alpha$) dari suatu populasi Normal sama dengan $\alpha_{0} $(konstanta) berdasarkan sampel random berukuran n. Langkah uji hipotesisnya dapat di urutkan sebagai berikut :\\
\begin{enumerate}
    \item Hipotesis
        \begin{enumerate}
            \item $H_{0} : \mu = \mu_{0}$ (uji dua sisi)\\
                $H_{1} : \mu \neq \mu_{0}$

            \item $H_{0} : \mu \leq \mu_{0}$ (uji sisi kanan)\\
                $H_{1} : \mu > \mu_{0}$

            \item $H_{0} : \mu \geq \mu_{0}$ (uji sisi kiri)\\
                $H_{1} : \mu < \mu_{0}$

        \end{enumerate}

    \item Diambil tingkat signifikansi $\alpha$

    \item Statistik penguji\\
        \begin{math}
            t = \frac{\bar{X} - \mu_{0}} {s/\sqrt{n}}
        \end{math}\\
        (jika $\sigma$ tidak diketahui) n atau p\_value

    \item Daerah kritis: daerah dimana $H_{0}$ ditolak.
        \begin{enumerate}
            \item $H_{0}$ ditolak jika $t > t_{n - 1;\alpha/2}$ atau $t < - t_{n-1;\alpha/2}$
            \item $H_{0}$ ditolak jika $t > t_{n - 1;\alpha/2}$ (untuk uji kanan) 
            \item $H_{0}$ ditolak jika $t < - t_{n-1;\alpha/2}$ (Untuk uji kiri) atau
            \item $H_{0}$ ditolak jika $p\_value < \alpha$ (untuk semua uji)
        \end{enumerate}
        
    \item Kesimpulan\\
        Berdasarkan langkah 4 dan hasil hitungan statistik penguji langkah 3, diambil kesimpulan apakah $H_{0}$ ditolak atau tidak ditolak pada tingkat signifikansi $\alpha$.
\end{enumerate}

\section{Pembahasan}
\subsection{Praktik}
\subsubsection{Praktik 1}
Berikut ini adalah data tekanan darah sistolik (dalam mmHg) 14 pasien yang menjalani terapi untuk hipertensi. Diasumsikan tekanan darah sistolik menyebar normal. Berdasarkan data berikut, dapatkah disimpulkan bahwa rata-rata tekanan darah pasien kurang dari 165 mmHg? Gunakan alpha 5\%\\
Berikut data tekanan darah pasien-pasien tersebut\\
\begin{center}
    183 152 178 157 194 163 144\\
    194 163 114 178 152 118 158
\end{center}

\paragraph{Jawab\\}


\end{document}
